\chapter{函数的概念与性质}

\section{函数的概念}

\begin{definition}[函数]
设 $A, B$ 是非空的实数集，如果对于集合 $A$ 中的任意一个数 $x$，按照某种确定的对应关系 $f$，在集合 $B$ 中都有唯一确定的数 $y$ 和它对应，那么就称 $f: A \to B$ 为从集合 $A$ 到集合 $B$ 的一个函数，记作：
\[ y = f(x), \quad x \in A \]
其中，$x$ 叫做自变量，$x$ 的取值范围 $A$ 叫做函数的定义域.
\end{definition}

\begin{exercise}
判断下列对应关系是否构成函数：$f(x) = \pm\sqrt{x}$ （$x \ge 0$）\selectbracket
\begin{enumerate}
    \item[A.] 构成
    \item[B.] 不构成
\end{enumerate}
\end{exercise}

\textbf{解析：}不构成。因为给定某一个非零的 $x$，会有两个对应的 $y$ 值（一正一负），这违反了函数“唯一确定”的定义。选 B。
